\section{Regional trajectories of cardiovascular mortality across the SSPs}

The Sellers (2020) regional decomposition — six geographic regions plus a separate High-Income bucket that pulls wealthy countries out of their geography — reveals three distinct trajectory shapes rather than one global pattern (Figure~\ref{fig:trajectory-deaths}). Regions that have already completed their cardiovascular epidemiological transition (Europe \& Central Asia) decline from the very start of the projection window under every SSP; regions in the middle of the transition (East Asia \& Pacific, High-Income Countries, Latin America \& Caribbean) peak mid-century before declining; and regions early in the transition (Sub-Saharan Africa, and under the weaker SSPs, the Middle East \& North Africa) do not peak within the horizon at all. Critically, the cardiovascular \emph{share} of all-cause deaths (Figure~\ref{fig:trajectory-share}) does not always track the same shape as absolute deaths — several regions show the share peaking years or decades before, or after, absolute deaths peak, a decoupling that is itself informative about the pace of each region's epidemiological transition under each SSP.

\begin{figure}[htbp]
  \centering
  \includegraphics[width=\textwidth]{Plots/RegionalSummaries/Sellers/fig5a_trajectory_deaths_by_region.png}
  \caption{Cardiovascular deaths (millions) by region, 1990--2100. Grey line: historical (1990--2023). Coloured lines: SSP1--SSP5 projections (2025--2100). Dashed vertical line marks the projection start. Sellers (2020) regional classification; y-axes are region-specific (free scale).}
  \label{fig:trajectory-deaths}
\end{figure}

\begin{figure}[htbp]
  \centering
  \includegraphics[width=\textwidth]{Plots/RegionalSummaries/Sellers/fig5b_trajectory_share_by_region.png}
  \caption{Cardiovascular deaths as a share of all-cause deaths by region, 1990--2100. Same colour scheme and layout as Figure~\ref{fig:trajectory-deaths}.}
  \label{fig:trajectory-share}
\end{figure}

\subsection*{East Asia \& Pacific}
Cardiovascular deaths rose steadily from 4.1 million in 1990 to 6.2 million by 2023 and continue rising under every SSP, peaking between 2045 (SSP5, 7.11 million) and 2055 (SSP3, 9.70 million) before declining to 2.60 million (15.1\% of all-cause deaths) under SSP5 versus 5.70 million (22.2\%) under SSP3 by 2100 — more than a two-fold gap in absolute deaths at century's end. The cardiovascular share itself is already past its historical peak (\textasciitilde41\% around 2017) and declines monotonically from 2025 under every SSP alike. This region's SSP1/SSP5 trajectories are consistent with the Healthy SSPs narrative's description of both pathways prioritizing comprehensive, technology-enabled care — SSP1 through ``highly comprehensive care, focused towards primary and preventive healthcare'' and SSP5 through digitized, AI-driven service delivery — while SSP3's higher and later peak matches its narrative of ``vulnerable infrastructure; overcrowded facilities and disrupted supply chains'' and care ``disproportionately focused on tertiary healthcare.''

\subsection*{Europe \& Central Asia}
This is the only region that has already passed its cardiovascular mortality peak historically (\textasciitilde2.7 million deaths circa 2003; 57\% cardiovascular share at its historical maximum) and declines monotonically across every SSP from 2025 onward. By 2100, deaths fall to 0.86 million (25\% share) under SSP1 versus 1.58 million (29\% share) under SSP3 — SSP3 slows the decline substantially but does not reverse it. This suggests that the region's pre-existing health-system infrastructure carries enough momentum that even the weakest-governance scenario (SSP3's ``fragmented, short-term planning with low accountability'') cannot arrest a decline already underway, though it can meaningfully blunt its pace.

\subsection*{High-Income Countries}
Unlike Europe \& Central Asia — with which it substantially overlaps geographically — the High-Income bucket was essentially flat historically (3.2--3.7 million, no strong secular trend 1990--2023), reflecting a broader mix of ageing high-income populations across regions (e.g., East Asian and North American high-income countries alongside European ones). Deaths still peak mid-century under every SSP (2045 for SSP1/SSP5, 2050 for SSP2/SSP3/SSP4) before declining to a 2100 range of 1.54 million (16.5\% share, SSP1) to 2.57 million (19.4\% share, SSP3). The SSP4 trajectory (2.05 million, 17.7\% share by 2100) sits close to SSP2, notably less divergent from SSP1/SSP5 than in most other regions — consistent with the narrative that SSP4's ``globally connected, technologically advanced elite'' benefits from ``high-performing institutions,'' a group this bucket disproportionately captures by construction, since the low-income populations that bear the brunt of SSP4's inequality are concentrated elsewhere.

\subsection*{Latin America \& Caribbean}
Deaths nearly doubled historically (0.53 to 0.99 million, 1990--2023) and continue rising under every SSP, peaking latest of the mid-transition regions: 2055 (SSP1/SSP5, 1.22--1.29 million) to 2070 (SSP3, 1.79 million), before declining to 0.64--1.46 million by 2100 (12.7\%--17.6\% share). Unlike East Asia, Europe, or High-Income Countries, the cardiovascular share does not peak immediately in 2025 but a few years into the projection window (2030--2035), indicating the region is still mid-transition rather than past it. This aligns with Table~1 of the Healthy SSPs narrative, where Latin America's urbanization and built-environment indicators sit closer to ``medium'' than the high-income regions, and with the narrative's description of uneven, urban-centric infrastructure development under SSP2--SSP4.

\subsection*{Middle East \& North Africa}
Already carrying one of the two highest cardiovascular shares of any region historically (\textasciitilde42--45\%, on par with Europe \& Central Asia), MENA shows the latest and least resolved peak of any region outside Sub-Saharan Africa. Under SSP3, deaths do not peak at all within the century — rising monotonically from 1.09 million (2025) to 2.27 million (2100, 27.8\% share), a 107\% increase. Even the best-case scenarios peak late (2070 under SSP1/SSP5). This is the region where the SSP3 narrative's combination of features — ``low funding, inadequate or non-existent for lower income regions,'' ``low access to hospitals or community clinics,'' and a ``strongly reactive'' governance posture — compounds an already elevated baseline cardiovascular burden rather than merely slowing an existing decline, as it does in Europe.

\subsection*{South Asia}
Deaths rose sharply from 1.67 to 4.04 million (1990--2023, +142\%). Under SSP1 and SSP5, 2025 is already the peak — these two scenarios flip the region into decline immediately, matching their pattern in East Asia \& Pacific — while SSP2--SSP4 continue rising to 2060--2075 before declining to 3.9--5.9 million by 2100. Notably, the cardiovascular \emph{share} peaks at 2025 under every SSP, including SSP2--SSP4, where absolute deaths keep climbing for decades afterward: population growth outpaces the growth in cardiovascular deaths even while the latter is still rising in absolute terms, so the share decouples from the raw death count earlier here than in any other region.

\subsection*{Sub-Saharan Africa}
This is the region furthest from completing its epidemiological transition, and the only one where the qualitative gap between SSPs is visible in the \emph{shape}, not just the level, of the trajectory. Deaths rose continuously from 0.58 million (1990) to 1.30 million (2023) with no sign of deceleration, and continue rising through 2100 under every SSP except SSP5 (which barely peaks, at 2095). SSP3 shows the largest proportional increase of any region-scenario combination in this analysis: 1.57 million (2025) to 6.48 million (2100), a 4.1-fold rise. The cardiovascular \emph{share} tells a more nuanced story than the death counts alone: under SSP1 and SSP2 — the scenarios with the strongest health-system investment — the share actually rises well above its 2025 level before falling back, peaking at 19.1\% around 2055 (SSP1) and 16.4\% around 2065 (SSP2), then declining to roughly 14\% by 2100. This bulge is consistent with an accelerated epidemiological transition: as SSP1/SSP2's investments in health systems, sanitation, and basic needs reduce mortality from infectious and other causes faster than cardiovascular mortality falls, cardiovascular disease temporarily becomes a larger share of a shrinking pool of deaths before eventually receding as cardiovascular-specific gains catch up. Under SSP3 and SSP4, by contrast, the share stays remarkably flat (12--13\%) across the entire century even as absolute deaths multiply several-fold — matching the narratives' description of persistent, multi-dimensional deprivation (``consistent, high'' poverty and hunger risk, degrading environmental management, weak health workforce) in which no single cause of death comes to dominate because other threats to survival remain elevated throughout, rather than cardiovascular disease ``taking over'' as the leading cause the way it does elsewhere.

\subsection*{Synthesis}
Across all seven regions, SSP3 consistently delays the cardiovascular mortality peak by 10--25 years relative to SSP1/SSP5 and produces a shallower post-peak decline, reproducing at the regional level the same SSP ordering (SSP3 worst, SSP1/SSP5 best) established globally. What the regional decomposition adds is where in the epidemiological transition each region sits: Europe \& Central Asia is already past its peak regardless of scenario; Sub-Saharan Africa and, under adverse governance, the Middle East \& North Africa are not projected to peak within the century at all. This maps directly onto the Healthy SSPs framework's core distinction between SSP1/SSP5 (``effective governance and adequate resourcing'') and SSP3/SSP4 (fragmented institutions, persistent inequity) — but shows that the practical consequence of that distinction differs by region: in already-transitioned regions it changes the \emph{pace} of an ongoing decline, while in early-transition regions it changes whether a decline happens within the century at all.
